\centerline{\bf \Huge Графы II}
\centerline{\bf \Huge Уровень: художка}

\begin{flushright}
        \scriptsize{}
\end{flushright}

\problem{1}
Докажите, что вершины триангуляции выпуклого n$-$угольника можно правильным образом покрасить в три цвета.

\problem{2}
Степени всех вершин графа не превосходят $d$. Докажите, что его вершины можно правильным образом раскрасить в $d^2+1$ цвет так, чтобы одноцветные вершины не имели общих соседей.

\problem{3}
Дан ориентированный граф, из каждой вершины которого выходит не более $d$ ребер. Докажите, что его вершины можно правильным образом раскрасить в $2 d+1$ цвет.

\problem{4}
В графе 1000 вершин, причем степень каждой не больше 9. Докажите, что можно выбрать такой подграф на 200 вершинах, что в нем не будет нечетных циклов.

\problem{5}
В графе любые два нечетных цикла имеют общую вершину. Докажите, что его вершины можно правильным образом покрасить в 5 цветов. Напомним, что покраску вершин графа называют правильной, если никакие две одноцветные вершины не соединены ребром.

\problem{6}
Некоторый граф правильно раскрашен в $k$ цветов, причём его нельзя правильно раскрасить в меньшее число цветов. Докажите, что в этом графе существует путь, вдоль которого встречаются вершины всех $k$ цветов ровно по одному разу.